\documentclass[12pt]{article}
\usepackage{geometry,fancyhdr,xr,hyperref,ifpdf,amsmath,rcs,indentfirst}
\usepackage{lastpage,longtable,Ventry,url,paunits,shortcuts,smallsec,color,tightlist,float}
\geometry{letterpaper,top=50pt,hmargin={20mm,20mm},headheight=15pt} 
\usepackage[stable]{footmisc}

\pagestyle{fancy} 

\RCS $Revision: 1.5 $
\RCS $Date: 2002/01/09 03:50:54 $

\fancypagestyle{first}{
\lhead{Comparing timescales for droplet growth and vapor diffusion}
\chead{}
\rhead{page~\thepage/\pageref{LastPage}}
\lfoot{} 
\cfoot{} 
\rfoot{}
}

\ifpdf
    \usepackage[pdftex]{graphicx} 
    \usepackage{hyperref}
    \pdfcompresslevel=0
    \DeclareGraphicsExtensions{.pdf,.jpg,.mps,.png}
\else
    \usepackage{hyperref}
    \usepackage[dvips]{graphicx}
    \DeclareGraphicsRule{.eps.gz}{eps}{.eps.bb}{`gzip -d #1}
    \DeclareGraphicsExtensions{.eps,.eps.gz}
\fi


\begin{document}
\newcommand{\vect}[1]{\boldsymbol{\vec{#1}}}
\pagestyle{first}


\section{Question}
\label{sec:question}



In the droplet growth notes we derived the growth rate
by making the assumption of
``steady droplet growth'', i.e.~

\begin{gather}
  \frac{ dm}{dt} = constant \\
\text{because: }\nonumber \\
\frac{ \partial \rho_v}{\partial t} = D \nabla^2 \rho_v = 0
\label{eq:steady}
\end{gather}

But how good is this assumption?  How does the time constant for
vapor changes compare to the time constant for droplet radius
changes? In symbols:

\begin{gather}
  \frac{1 }{\rho_v} \frac{ d \rho_v}{dt} = \frac{ 1}{\tau_\rho} \\
\text{compare to the time constant: }\nonumber \\
\frac{1}{r} \frac{ dr}{dt} = \frac{ 1}{\tau_r} 
\end{gather}


We showed in class that the droplet time scale $\tau^{-1}_r = \frac{ 1}{r}\frac{ dr}{dt}   \sim 1\ second$.  
To find $\tau_\rho$ the timescale for change in  vapor density $\rho_v$, you need
the diffusion coefficient $D$, which according to constants.m is $D \approx 2.35 \times 10^{-5}
\ \un{m^2\,s^{-1}}$.  Do a scale analysis for the neighborhood within  10 radii of  a 10 \mum drop and
show that 


\begin{gather}
  \frac{1 }{\rho_v} \frac{ d \rho_v}{dt} =  D \frac{1 }{\rho_v} \nabla^2 \rho_v = D  \frac{1 }{\rho_v} \frac{ d^2 \rho_v}{dr^2} 
\end{gather}
does in fact give a timescale much less than 1 second.

\section{Solution}
\label{sec:solution}


We need to make a guess at the size of $\nabla^2 \rho_v \sim \frac{ d}{dx} \left ( \frac{ d \rho_v}{dx} \right )$.  If we were
going to evaluate this with matlab, it would look something like

\begin{gather}
  \nabla^2 \rho_v \sim \frac{ d}{dx} \left( \frac{ d \rho_v}{dx} \right ) = \\
\frac{\frac{\rho_{v2} - \rho_{v1}}{\Delta x} - \frac{\rho_{v1} - \rho_{v0}}{\Delta x}}{\Delta x}
= \frac{(\rho_{v2} + \rho_{v0}) - 2 \rho_{v1}}{\Delta x^2} = \frac{2 \overline{\rho_{v}} - 2 \rho_{v1}}{\Delta x^2}   \sim \frac{2\Delta \rho_v }{(\Delta x)^2}  \\
\text{Make a rough guess: }\nonumber \\
2\Delta \rho_v \sim \rho_v\ \text{(upper limit)}\\
\Delta x \sim 100\ \mum \ \text{(10 radii away from surface)}\nonumber 
\end{gather}




Which means that:

\begin{gather}
 \frac{ 1}{\tau_\rho} = \frac{1 }{\rho_v} \frac{ d \rho_v}{dt} =  \frac{D }{\rho_v}  \nabla^2 \rho_v 
\sim \frac{ D \rho_v}{\rho_v (\Delta x)^2} \sim \frac{10^{-5}}{(10^{-4})^2} \sim 1000\ \un{s^{-1}}\\
\text{or in other words} \nonumber \\
\tau_\rho \sim 0.001\ \un{s^{-1}} \ll 1\ \text{second}
\end{gather}
This means that the vapor field adjusts very quickly to changes compared to a 5 \mum cloud droplet,
So ``steady growth'' is a good assumption for atmospheric water vapor.




\end{document}


%%% Local Variables:
%%% mode: latex
%%% TeX-master: t
%%% End:
